\chapter{The maximum principle}
\label{chap:maximum-principle}
\label{secmaxprin}

The scalar, tensor, strong, and Harnack maximum principles below are used
for Ricci flow; the tensor maximum principle is due to Hamilton
\cite{Hamilton4MPIC}. The tensor results also use the formulations in
\cite{Hamilton4MPCO} and \cite{ChowKnopf}.

\section{Maximum principle for scalar curvature}

\begin{proposition}[Evolution bound for minimum scalar curvature]
\label{prop:scalar-curvature-min-evolution}
\lean{MorganTianLib.monotoneOn_scalarCurvatureMinimum_of_isRicciFlowOn_Ico, MorganTianLib.scalarCurvatureMinimum_nonnegative_lower_bound_of_isRicciFlowOn, MorganTianLib.scalarCurvatureMinimum_negative_lower_bound_of_isRicciFlowOn, MorganTianLib.lifespan_le_of_isRicciFlowOn_of_scalarCurvatureMinimum_pos}
\leanok
\uses{def:ricci-flow, def:ricci-curvature}
\label{scalarevol}
Let $(M,g(t)),\ 0\le t<T$, be a Ricci flow on a compact
$n$-manifold, and let $R_\mathrm{min}(t)$ be the minimum scalar curvature.
Then $R_\mathrm{min}$ is non-decreasing. If $R_\mathrm{min}(0)\ge0$, then
\[
R_\mathrm{min}(t)\ge R_\mathrm{min}(0)\left(\frac{1}{1-\frac{2t}{n}R_\mathrm{min}(0)}\right).
\]
If $R_\mathrm{min}(0)>0$, then $T\le n/(2R_\mathrm{min}(0))$.
If $R_\mathrm{min}(0)<0$, then
\[
R_\mathrm{min}(t)\ge-\frac{n|R_\mathrm{min}(0)|}{2t|R_\mathrm{min}(0)|+n}.
\]
\end{proposition}
\begin{proof}
  \uses{lem:forward-difference-quotient, thm:curvature-evolution, def:ricci-curvature}

\sketch The evolution equation is $\partial_tR=\Delta R+2|\Ric|^2$.
At a minimum, $\Delta R\ge0$ and $|R|^2\le n|\Ric|^2$, giving
the two claims below. Integrating the resulting forward-difference inequality
for $R_\mathrm{min}$ gives monotonicity and the stated bounds. For
$R_\mathrm{min}$ nonzero throughout $[0,t]$, use
\[
S(t)=-R_\mathrm{min}(t)^{-1}-\frac{2t}{n}+R_\mathrm{min}(0)^{-1},
\qquad S(0)=0,\quad S'(t)\ge0,
\]
and hence
\begin{equation}\label{Rminineq}
\frac{1}{R_\mathrm{min}(t)}\le\frac{1}{R_\mathrm{min}(0)}-\frac{2t}{n}.
\end{equation}
\end{proof}

\begin{claim}[Time derivative of R at its minimum]
\label{claim:scalar-min-derivative-bound}
\uses{def:ricci-curvature}
\lean{MorganTianLib.scalar_min_deriv_bound_of_isRicciFlowOn}
\label{firstclaim}
If $R(x,t)=R_\mathrm{min}(t)$, then
$(\partial R/\partial t)(x,t)\ge\frac{2}{n}R^2(x,t)$.
\end{claim}
\begin{proof}
\sketch At a minimum $\Delta R\ge0$; combine the scalar-curvature evolution
with $|R|^2\le n|\Ric|^2$.
\end{proof}

\begin{claim}[Forward-difference bound for the minimum scalar curvature]
\label{claim:scalar-min-forward-difference}
\uses{claim:scalar-min-derivative-bound}
\lean{MorganTianLib.forwardDiffQuotientGE_scalarCurvatureMinimum_of_isRicciFlowOn_Ico}
\leanok
\label{Rminprime}
\[
\frac{d}{dt}R_\mathrm{min}(t)\ge\frac{2}{n}R_\mathrm{min}^2(t),
\]
with forward difference quotients at non-smooth times.
\end{claim}
\begin{proof}
\uses{prop:forward-difference-maximum}
\sketch Apply the forward-difference maximum proposition to the preceding
pointwise minimum estimate.
\end{proof}

\section{The maximum principle for tensors}

For a closed convex set $Z$ in a finite-dimensional real vector space $V$,
let $T_zZ$ be the tangent cone at its boundary and set $T_zZ=V$ in the
interior. Thus $v\notin T_zZ$ exactly when an affine linear functional
vanishing at $z$ is non-positive on $Z$ and positive on $v$.

\begin{definition}[Vector field preserving a convex set]
\label{def:vector-field-preserves-convex-set}
For a closed convex $Z\subset V$, a smooth vector field $\psi$ on a
neighborhood of $Z$ preserves $Z$ if $\psi(z)\in T_zZ$ for every $z\in Z$.
\end{definition}

\begin{lemma}[Integral-curve criterion for preserving a convex set]
\label{lem:integral-curve-preserves-convex-set}
\uses{def:vector-field-preserves-convex-set}

For closed convex $Z\subset V$ and smooth $\psi$ near $Z$, $\psi$ preserves
$Z$ iff every integral curve of $\psi$ starting in $Z$ remains in $Z$.
This is the standard convex-set criterion \cite{Hamilton4MPCO}.
\end{lemma}

\subsection{The global version}\label{4.3.3}

\begin{definition}[Convex subbundle and fiberwise preserving vector fields]
\label{def:convex-subbundle-preserving-vector-field}
\uses{def:vector-field-preserves-convex-set}

For a vector bundle $\pi:{\mathcal V}\to M$ and closed ${\mathcal Z}\subset
{\mathcal V}$, call ${\mathcal Z}$ convex when every fiber $Z_x$ is convex.
A fiberwise vector field on a neighborhood ${\mathcal U}$ preserves ${\mathcal Z}$
when its restriction to each fiber preserves $Z_x$.
\end{definition}

\begin{theorem}[Maximum principle for tensors (global version)]
\label{thm:maximum-principle-tensors-global}
\uses{def:convex-subbundle-preserving-vector-field, thm:levi-civita-connection}
\label{MP}
Let $(M,g)$ be compact, ${\mathcal V}\to M$ a tensor bundle, and ${\mathcal Z}$
a closed convex parallel-invariant subset. If a fiberwise vector field $\psi$
preserves ${\mathcal Z}$ and
\[
\partial_t{\mathcal T}=\Delta{\mathcal T}+\psi({\mathcal T}),
\]
then ${\mathcal T}(x,0)\in{\mathcal Z}$ for all $x$ implies
${\mathcal T}(x,t)\in{\mathcal Z}$ for all $x,t$.
The corresponding proof is given in \cite{Hamilton4MPCO}, \cite{Hamilton4MPCO},
and \cite{Hamilton4MPCO}.
\end{theorem}

\begin{theorem}[Maximum principle for tensors with time-varying convex sets]
\label{thm:maximum-principle-tensors-varying-z}
\uses{def:convex-subbundle-preserving-vector-field, thm:levi-civita-connection}
\label{varyingZ}
Let ${\mathcal Z}\subset{\mathcal V}\times[0,T]$ have closed, convex,
parallel-invariant time slices. If $\psi$ preserves the family in the
integral-curve sense and $\partial_t{\mathcal T}=\Delta{\mathcal T}+\psi({\mathcal T})$,
then ${\mathcal T}(x,0)\in{\mathcal Z}(0)$ implies
${\mathcal T}(x,t)\in{\mathcal Z}(t)$.
The time-dependent formulation is also recorded in \cite{Hamilton4MPCO}.
\end{theorem}

\subsection{The local version}

\begin{theorem}[Local maximum principle for tensors]
\label{thm:maximum-principle-tensors-local}
\uses{def:convex-subbundle-preserving-vector-field, thm:levi-civita-connection}
\label{localMP}
Let $\overline U\subset M$ be compact, smooth, connected, and codimension
zero, and let ${\mathcal Z}\subset{\mathcal V}$ be closed, convex, and
parallel-invariant. If $\psi$ preserves ${\mathcal Z}$ and
${\mathcal T}$ satisfies $\partial_t{\mathcal T}=\Delta{\mathcal T}+\psi({\mathcal T})$,
then ${\mathcal T}(x,0)\in{\mathcal Z}$ on $\overline U$ and, if boundary
values remain in ${\mathcal Z}$, then ${\mathcal T}(x,t)\in{\mathcal Z}$ on
\overline U\times[0,T]$.
This local form is proved in \cite{Hamilton4MPCO}.
\end{theorem}

\section{Applications of the maximum principle}
The curvature evolution formulas are used in the standard form
\cite{ChowLuNi} and \cite{ChowLuNi}; the three-dimensional roundness result is
\cite{Hamilton3MPRC} and \cite{Hamilton3MPRC}.
\subsection{Ricci flows with normalized initial conditions}

\begin{definition}[Normalized initial conditions for Ricci flow]
\label{def:normalized-initial-conditions}
\lean{MorganTianLib.IsNormalizedInitialConditions}
\leanok
\uses{def:ricci-flow, def:riemann-curvature-tensor}
\label{norminitcond}
A Ricci flow has normalized initial conditions when its initial time is zero,
$|\Rm(x,0)|\le1$, and, with $\omega_n$ the unit Euclidean $n$-ball volume,
\[
\Vol(B(x,0,r))\ge(\omega_n/2)r^n\qquad (r\le1).
\]
\end{definition}

\begin{proposition}[Short-time curvature and volume bounds for normalized flows]
\label{prop:normalized-flow-short-time-bounds}
\uses{def:normalized-initial-conditions}
\label{kappa0r0t0}
There is $\kappa_0>0$ depending only on $n$ such that a complete bounded-
curvature normalized flow satisfies $|\Rm(x,t)|\le2$ and
$\Vol B(x,t,r)\ge\kappa_0r^n$ for $0\le t\le\min(T,2^{-4})$ and $r\le1$.
\end{proposition}
\begin{proof}
\uses{thm:curvature-evolution}
\sketch Curvature evolution gives the bound on $[0,2^{-4}]$. Distance
distortion and the volume evolution equation, together with the initial
volume bound, give $B(x,0,r_0r)\subset B(x,t,r)$ and
$\Vol_t(B(x,0,s))\ge A^{-1}\Vol_0(B(x,0,s))$, hence
$\Vol_t(B(x,t,r))\ge A^{-1}(\omega_n/2)r_0^nr^n$.
\end{proof}

\subsection{Extending flows}
\begin{proposition}[Extension criterion for Ricci flows]
\label{prop:ricci-flow-extension-criterion}
\uses{def:ricci-flow}
\label{flowextend}
For a compact Ricci flow on $[0,T)$ with $T<\infty$, either the flow extends
past $T$ or $|\Rm|$ is unbounded on $M\times[0,T)$.
\end{proposition}

\subsection{Non-negative curvature is preserved}

\begin{lemma}[Curvature-operator evolution ODE in dimension three]
\label{lem:curvature-operator-ode-dimension-three}
\uses{prop:curvature-operator-evolution, def:curvature-operator}
\label{quad}
For a three-dimensional curvature operator ${\mathcal T}$ in an evolving
frame,
\[
\partial_t{\mathcal T}=\triangle{\mathcal T}+\psi({\mathcal T}),
\qquad \psi({\mathcal T})={\mathcal T}^2+{\mathcal T}^{\sharp}.
\]
If ${\mathcal T}=\operatorname{diag}(\lambda,\mu,\nu)$, then
\[
\psi({\mathcal T})=\operatorname{diag}(\lambda^2+\mu\nu,
\mu^2+\lambda\nu,\nu^2+\lambda\mu).
\]
\end{lemma}

\begin{corollary}[Non-negative curvature operator is preserved]
\label{cor:nonnegative-curvature-operator-preserved}
\uses{def:curvature-operator}
\label{rm>0}
For a compact connected three-dimensional Ricci flow, $\Rm(x,0)\ge0$ implies
\Rm(x,t)\ge0$ for all $x,t$.
\end{corollary}
\begin{proof}
\uses{thm:maximum-principle-tensors-global,lem:curvature-operator-ode-dimension-three}
\sketch The smallest-eigenvalue superlevel set is closed, convex, and
parallel-invariant. The displayed ODE has nonnegative vector-field direction
on this set, so apply the global tensor maximum principle.
\end{proof}

\begin{corollary}[Non-negative Ricci curvature is preserved]
\label{cor:nonnegative-ricci-preserved}
\uses{def:ricci-curvature}
\label{Ric>0}
For a compact connected three-dimensional Ricci flow, $\Ric(x,0)\ge0$ implies
$\Ric(x,t)\ge0$ for all $x,t$.
\end{corollary}
\begin{proof}
\uses{thm:maximum-principle-tensors-global,lem:curvature-operator-ode-dimension-three}
\sketch Apply the tensor maximum principle to the convex cone of curvature
operators whose sum of the two smallest eigenvalues is nonnegative.
\end{proof}

\section{The strong maximum principle for curvature}
\begin{theorem}[Strong maximum principle for the heat equation]
\label{thm:strong-maximum-principle-heat}\label{SMPheat}
For a compact connected manifold with boundary, a Dirichlet heat solution
$h$ with $h(\cdot,0)\ge0$ remains nonnegative; if it is positive at one
initial point, it is positive on the interior for positive times.
\end{theorem}

\begin{proposition}[Strong maximum principle for curvature tensors]
\label{prop:strong-maximum-principle-curvature}
\uses{def:convex-subbundle-preserving-vector-field}
\label{SMPcurv}
Let $s$ be fiberwise convex and parallel-invariant, and let
$\partial_t{\mathcal T}=\Delta{\mathcal T}+\psi({\mathcal T})$ on a compact
codimension-zero domain. Assume $\psi(A)$ lies in the tangent cone of
$\{B:s(B)\ge s(A)\}$ whenever $s(A)\ge0$. If $s({\mathcal T}(x,0))\ge0$,
the boundary remains nonnegative, and $s({\mathcal T}(x_0,0))>0$ at an
interior point, then $s({\mathcal T}(x,t))>0$ in the interior for $t>0$.
\end{proposition}
\begin{proof}
\uses{thm:strong-maximum-principle-heat,thm:maximum-principle-tensors-local}
\sketch Choose a nonnegative Dirichlet heat solution $h$ below
$s({\mathcal T}(\cdot,0))$, positive at $x_0$, and apply the local tensor
maximum principle to the convex set $s({\mathcal T})\ge h\ge0$.
\end{proof}

\begin{theorem}[Vanishing scalar curvature forces a flat metric]
\label{thm:vanishing-scalar-curvature-forces-flat}\label{firststrongmax}
For a connected three-dimensional Ricci flow with nonnegative sectional
curvature, $R(p,T)=0$ implies that the flow is flat on $[0,T]$.
\end{theorem}
\begin{proof}
\uses{thm:curvature-evolution}
\sketch If curvature is positive at an earlier point, compare $R$ with a
positive Dirichlet heat solution using
$\partial_t(R-h)=\triangle(R-h)+2|\Ric|^2$; the strong maximum principle
forces $R(p,T)>0$, a contradiction.
\end{proof}

\begin{corollary}[Local product splitting at a Ricci null eigenvalue]
\label{cor:local-product-splitting-ricci-null-eigenvalue}
\uses{def:curvature-operator}
\label{localprod}
For a nonflat connected three-dimensional flow with nonnegative sectional
curvature, a zero Ricci eigenvalue at time $T>0$ implies that for every
$t\in(0,T]$ the flow locally splits as a positively curved surface times a
line, with the product flow structure.
\end{corollary}
\begin{proof}
\uses{cor:nonnegative-ricci-preserved,thm:curvature-evolution,prop:strong-maximum-principle-curvature,thm:vanishing-scalar-curvature-forces-flat}
\sketch Apply the strong principle to the convex function given by the sum
of the two smallest curvature-operator eigenvalues. It vanishes everywhere,
so the curvature operator has eigenvalues $\{\lambda,0,0\}$ with $\lambda>0$.
The null line field is parallel and time-invariant; its orthogonal complement
therefore gives the asserted local product.
\end{proof}

\begin{corollary}[Global product splitting on a cover]
\label{cor:global-product-splitting-cover}
\uses{cor:local-product-splitting-ricci-null-eigenvalue}
\label{nullspace}
For a complete connected nonflat three-dimensional flow with $\Rm\ge0$ and a
zero curvature-operator eigenvalue at time $T$, a cover splits as
$(N,h(t))\times(\Ar,ds^2)$, where $N$ has positive curvature and the flow is
the product with the trivial line flow.
\end{corollary}
\begin{remark}[Classification of covers for the product splitting]
\label{rem:product-splitting-cover-classification}
\uses{cor:global-product-splitting-cover}

The possible cover is trivial, a normal $\Zee$-cover, a two-sheeted cover, or
a normal infinite-dihedral cover; the last two carry only a non-orientable
null line field.
\end{remark}

\begin{proposition}[Non-flat cones do not arise from non-negatively curved Ricci flow]
\label{prop:no-nonflat-cones-from-ricci-flow}
\uses{def:cone}
\label{nocones}
An open subset of a three-dimensional nonnegatively curved Ricci flow that is
isometric at positive time to an open subset of a cone is flat for all times.
\end{proposition}
\begin{proof}
\uses{prop:cone-curvature,thm:vanishing-scalar-curvature-forces-flat,cor:local-product-splitting-ricci-null-eigenvalue}
\sketch A nonflat cone has a two-dimensional curvature null space and a
nonconstant nonzero eigenvalue along cone lines, whereas local splitting makes
that eigenvalue constant along the parallel null line. Contradiction.
\end{proof}

\begin{theorem}[Compact positive-Ricci flows become round; corrected form]
\label{thm:nonnegative-ricci-flow-becomes-round}
\uses{def:curvature-operator,cor:nonnegative-ricci-preserved}
\label{flowtoround}
Let $(M,g(t))$, $0\le t<T$ with $T>0$, be a maximal Ricci flow on a nonempty compact
connected three-manifold without boundary, with $\Ric(g(0))$ positive
definite. Then $T<\infty$ and $\Ric(g(t))$ remains positive definite.
Let $\lambda(x,t)$ and $\nu(x,t)$ be the largest and smallest eigenvalues
of the self-adjoint curvature operator on $\wedge^2T_xM$, using the induced
exterior metric and the normalization
$\Rm(u\wedge v,u\wedge v)=K(u\wedge v)$ for orthonormal $u,v$. Then
\[
\lim_{t\nearrow T}\frac{\max_{x\in M}\lambda(x,t)}{\min_{x\in M}\nu(x,t)}=1,
\]
and $\lambda(x,t)\to\infty$ for every $x\in M$. For any fixed $p\in M$,
the spatially constant rescaling $\lambda(p,t)g(t)$ converges smoothly on
$M$ as $t\nearrow T$ to a metric of constant sectional curvature $1$.
Consequently, $M$ is diffeomorphic to a three-dimensional spherical
space-form. This is the positive-Ricci form of Hamilton's theorem
\cite{Hamilton3MPRC}; the printed nonnegative-Ricci dichotomy is excluded
by the shrinking product flow on $S^2\times S^1$.
\end{theorem}

\subsection{Solitons of positive curvature}
\begin{theorem}[Positive Ricci solitons are round spherical space forms]
\label{thm:positive-ricci-soliton-round}\label{sphsf}
A compact three-dimensional soliton with positive Ricci curvature is round,
and is a quotient of the round $S^3$ by a finite freely acting subgroup of
$O(4)$.
\end{theorem}
\begin{proof}
\uses{thm:nonnegative-ricci-flow-becomes-round}
\sketch The maximal flow becomes round at its finite singular time by
Theorem~\ref{flowtoround}; soliton self-similarity then makes every time slice
round.
\end{proof}
\begin{remark}[Extending roundness to shrinking solitons]
\label{rem:roundness-shrinking-soliton-preview}
\uses{thm:positive-ricci-soliton-round, def:gradient-shrinking-soliton}

The pinching theorem below extends the roundness conclusion to compact
shrinking solitons without assuming positive Ricci curvature.
\end{remark}

\section{Pinching toward positive curvature}
The pinching argument follows the curvature maximum-principle estimates in
\cite{HamiltonNSRF3M} and \cite{HamiltonNSRF3M}; the three-dimensional soliton consequence is due to
Ivey \cite{Ivey}.
\begin{theorem}[Pinching toward positive curvature]
\label{thm:pinching-toward-positive-curvature}
\uses{def:ricci-flow, def:curvature-operator}
\label{pinch}
For a compact three-dimensional flow whose initial curvature-operator
eigenvalues satisfy $\lambda(x,0)\ge\mu(x,0)\ge\nu(x,0)\ge-1$, set
$X=\max(-\nu,0)$. Then
\begin{enumerate}
\item $R(x,t)\ge-6/(4t+1)$;
\item if $X(x,t)>0$, then
\[
R(x,t)\ge2X(x,t)(\log X(x,t)+\log(1+t)-3).
\]
\end{enumerate}
\end{theorem}
\begin{proof}
\uses{prop:scalar-curvature-min-evolution,thm:maximum-principle-tensors-varying-z}
\sketch The scalar estimate gives the first pinching inequality. For the
second, apply the time-varying tensor maximum principle to the three helper
claims below. Since $S=R/2$, the resulting inequalities imply the stated bound
for $X\ge1/(1+t)$; for $0<X\le1/(1+t)$ use $S\ge-3X$ and
$f_t(X)<-3X$.
\end{proof}
\begin{remark}[Consequence: negative eigenvalues are small relative to scalar curvature]
\label{rem:pinching-negative-eigenvalues-small}
\uses{thm:pinching-toward-positive-curvature}

As scalar curvature tends to infinity, the absolute values of negative
curvature-operator eigenvalues are negligible relative to it.
\end{remark}

\begin{claim}[Convexity of the pinching region]
\label{claim:pinching-region-convex}
For each $x,t$, the fiber $Z(x,t)$ of the pinching region is a convex subset
of $\Sym^2(\wedge^2T_x^*M)$.
\end{claim}
\begin{proof}
\sketch Let $f_t(x)=x(\log x+\log(1+t)-3)$ and let $\xi(t)>1/(1+t)$ be the
second solution of $f_t(x)=-3/(1+t)$. On $[\xi(t),\infty)$, $f_t'>0$
and $f_t''>0$, so the associated $(S,X)$ region is convex. Since the
smallest eigenvalue is convex, $X=\max(-\nu,0)$ is concave, which proves
convexity of the inverse image.
\end{proof}

\begin{claim}[Initial curvature lies in the pinching region]
\label{claim:pinching-initial-condition}
\uses{claim:pinching-region-convex}

${\mathcal T}(x,0)\in Z(x,0)$ for all $x\in M$.
\end{claim}
\begin{proof}
\sketch The initial lower bound $\lambda+\mu+\nu\ge-3$ gives
$S\ge-3$. If $0<X(x,0)\le1$, then $S\ge-3X\ge X(\log X-3)$.
\end{proof}

\begin{claim}[Non-negativity of the auxiliary quantity I]
\label{claim:pinching-auxiliary-quantity-nonnegative}
\lean{MorganTianLib.pinching_auxiliary_nonneg_of_ordered_eigenvalues}
\leanok
$I\ge0$, where $X=-\nu$, $Y=-\mu$, and
\[
I=XY^2+\lambda Y(Y-X)+\lambda^2(X-Y).
\]
\end{claim}
\begin{proof}
\sketch If $Y\le0$, then $\lambda\ge0$ and $X\ge Y$, so $I\ge0$. If
$Y>0$, then
\[
I=Y^3+(X-Y)(\lambda^2-\lambda Y+Y^2)>0,
\qquad \lambda^2-\lambda Y+Y^2=(\lambda-Y/2)^2+3Y^2/4.
\]
\end{proof}

\begin{claim}[The vector field preserves the pinching family of convex sets]
\label{claim:pinching-vector-field-preserves-family}
\uses{claim:pinching-region-convex}

The vector field $\psi({\mathcal T})={\mathcal T}^2+{\mathcal T}^{\sharp}$
preserves ${\mathcal Z}(t)$.
\end{claim}
\begin{proof}
\uses{claim:pinching-auxiliary-quantity-nonnegative}
\sketch Write $S=\operatorname{tr}{\mathcal T}$ and $X=\max(-\nu,0)$.
Along an integral curve,
\begin{equation}\label{Sineq}S(t)\ge-3/(1+t).\end{equation}
For $X>0$, $Y=-\mu$ gives
\begin{align*}
\dot X&=-X^2+Y\lambda,\\
\dot S&=X^2+Y^2+\lambda^2+XY-\lambda(X+Y).
\end{align*}
Moreover,
\begin{equation}\label{dX}X\dot S-(S+X)\dot X=X^3+I.
\end{equation}
The auxiliary claim implies
\begin{equation}\label{dX1}X\dot S-(S+X)\dot X\ge X^3.\end{equation}
For $W=S/X-\log X$,
\begin{equation}\label{dW}\dot W\ge X.
\end{equation}
At a first boundary contact, the $S=-3/(1+t)$ boundary is excluded by
\eqref{Sineq}; on $S=f_t(X)$, $X>1/(1+t)$ and the inequality for $W$
excludes crossing. Thus the family is preserved.
\end{proof}

\begin{theorem}[Pinching toward positive curvature persists from time a]
\label{thm:pinching-persists-from-time-a}
\uses{thm:pinching-toward-positive-curvature}
\label{pincha}
For $a\ge0$, assume the two pinching inequalities at time $a$ for a compact
three-dimensional flow. Then for $a\le t<T$,
\begin{eqnarray}\label{Rlower}
R(x,t)&\ge&-6/(4t+1)\\\label{RXineq}
R(x,t)&\ge&2X(x,t)(\log X(x,t)+\log(1+t)-3),
\end{eqnarray}
whenever $X(x,t)>0$.
\end{theorem}

\begin{corollary}[Curvature bound from a scalar-curvature bound under pinching]
\label{cor:curvature-bound-from-scalar-pinching}
\uses{thm:pinching-persists-from-time-a}

Under the hypotheses of the preceding theorem there is a continuous $\phi$
such that $R(x,t)\le R_0$ implies $|\Rm(x,t)|\le\phi(R_0)$.
\end{corollary}
\begin{proof}
\uses{thm:pinching-persists-from-time-a}
\sketch If $X=0$, $|\Rm|\le R/2$; if $X>0$, pinching bounds $X$ by $R_0$
for large $R_0$, and hence bounds $\lambda$ and $\nu$, which bounds $|\Rm|$.
\end{proof}

\begin{definition}[Curvature pinched toward positive for a generalized Ricci flow]
\label{def:curvature-pinched-toward-positive}
\uses{def:generalized-ricci-flow, thm:pinching-toward-positive-curvature, def:curvature-operator}
\label{pinchd}
A generalized Ricci flow on $[0,\infty)$ is curvature pinched toward positive
when, for every $x$,
\begin{enumerate}
\item $R(x)\ge-6/(4{\bf t}(x)+1)$;
\item $R(x)\ge2X(x)(\log X(x)+\log(1+{\bf t}(x))-3)$ whenever $X(x)>0$,
\end{enumerate}
where $X$ is the maximum of zero and the negative smallest eigenvalue of
$\Rm$.
\end{definition}

\begin{theorem}[Compact shrinking solitons are round]
\label{thm:shrinking-soliton-round}
A compact three-dimensional shrinking soliton is round.
\end{theorem}
\begin{proof}
\uses{thm:nonnegative-ricci-flow-becomes-round,thm:pinching-toward-positive-curvature,cor:nonnegative-curvature-operator-preserved}
\sketch Rescale so the initial curvature-operator eigenvalues have absolute
value at most one. If $X(x,0)>0$, soliton scaling makes $X/R$ fixed while
the pinching estimate forces it to vanish as the singular time is approached;
hence $\Rm\ge0$. The preservation and roundness theorems then exclude the
flat case and imply roundness of every soliton slice.
\end{proof}

\begin{theorem}[Compact three-dimensional Ricci solitons are Einstein]
\label{thm:ricci-soliton-einstein-dimension-three}
\uses{def:einstein-manifold}
\label{ivey}
Any compact three-dimensional Ricci soliton $g_0$ is Einstein.
\end{theorem}

\subsection{The Harnack inequality}
The differential and integrated inequalities are Hamilton's results; the
general soliton statement is recorded in \cite{Ivey}.
\begin{theorem}[Hamilton's Harnack inequality for Ricci flow]
\label{thm:hamilton-harnack-inequality}
\uses{def:ricci-curvature, def:curvature-operator}
\label{Harnack}
For a complete Ricci flow with bounded nonnegative curvature operator on
$(T_0,T_1)$ and any time-dependent vector field $\chi$,
\[
\partial_tR+\frac{R}{t-T_0}+2\langle\chi,\nabla R\rangle+2\Ric(\chi,\chi)\ge0,
\qquad \partial_tR+\frac{R}{t-T_0}\ge0.
\]
This is Hamilton's differential Harnack inequality \cite{Hamiltonharnack} and
\cite{Hamiltonharnack}.
\end{theorem}
\begin{remark}[Harnack inequality special case with zero vector field]
\label{rem:harnack-inequality-zero-field-case}
\uses{thm:hamilton-harnack-inequality}

The second inequality is the case $\chi=0$.
\end{remark}

\begin{corollary}[Scalar curvature is non-decreasing for ancient flows]
\label{cor:ancient-flow-scalar-curvature-nondecreasing}
\uses{thm:hamilton-harnack-inequality}
\label{posderiv}
For an ancient complete flow with bounded nonnegative curvature operator,
$\partial_tR(x,t)\ge0$.
\end{corollary}
\begin{proof}
\uses{thm:hamilton-harnack-inequality}
\sketch Set $\chi=0$ and let $T_0\to-\infty$ in the differential Harnack
inequality.
\end{proof}

\begin{theorem}[Integrated Harnack inequality]
\label{thm:harnack-inequality-integrated}
\uses{def:geodesic}
\label{Harnack2}
For a complete flow with bounded nonnegative curvature operator on
$[t_1,t_2]$ and points $x_1,x_2$,
\[
\log\left(\frac{R(x_2,t_2)}{R(x_1,t_1)}\right)
\ge-\frac12\frac{d_{t_1}^2(x_2,x_1)}{t_2-t_1}.
\]
\end{theorem}
\begin{proof}
\uses{thm:hamilton-harnack-inequality}
\sketch Take $\chi=-\nabla\log R/2$ in the differential Harnack inequality and
divide by $R$ to obtain
$\partial_t\log R-|\nabla\log R|^2/2\ge0$. Integrate along the constant-speed
$g(t_1)$-geodesic from $x_1$ to $x_2$; completing the square and using
$\Ric\ge0$ yields the stated distance-energy bound.
\end{proof}
